\section{Factorization, arithmetic clocks and hierarchy of traces}
\label{lectura:manuscrito-02}
The native product distinguishes two properties: irreducibility of a normal form and the period of a multiplicative action on it. This chapter constructs their spectral readouts, preserves the normalizations of each functional and prepares the completed form developed in Section~\ref{lectura:manuscrito-03}.

\subsection{From the irreducible to its clock}
Product and coproduct determine which objects are irreducible. The arithmetic clock describes another property of those objects: return of a finite multiplicative orbit. For an integer $n\ge2$ coprime to ten, define

\[
\tau_n=\min\{k\ge1:10^k\equiv1\pmod n\}.
\]
The length $\tau_n$ corresponds to the orbit of the unit under multiplication by ten in the group of units modulo $n$. The decimal base fixes this readout; another coprime base determines another clock. Irreducibility, by contrast, belongs to the product constructed before this periodic representation is chosen.

\HMTresultado{Proposition}{Synchronization of clocks}{rz-num-02-1}
If $n=\prod_p p^{a_p}$ and $\gcd(n,10)=1$, then

\[
\tau_n=\operatorname{mcm}_{p^{a_p}\parallel n}\tau_{p^{a_p}}.
\]
\textbf{Proof.} An exponent $k$ satisfies $10^k\equiv1\pmod n$ if and only if it does so modulo each prime power in the factorization. This is equivalent to $k$ being a multiple of every order $\tau_{p^{a_p}}$. The smallest common positive exponent is the least common multiple. \hfill\(\square\)

Factorization therefore organizes synchronization. Powers of the same prime must be retained as such: their period is not obtained by replacing them, without proof, with the base prime's period.

\HMTresultado{Proposition}{Nonadic residual closure}{rz-num-02-2}
If $\gcd(n,90)=1$, the integer block

\[
M_n=\frac{10^{\tau_n}-1}{n}
\]
is divisible by nine.

\textbf{Proof.} The period definition makes $M_n$ integral. Since $10\equiv1\pmod9$, the numerator is divisible by nine. The integer $n$ is invertible modulo nine, so $nM_n\equiv0\pmod9$ implies $M_n\equiv0\pmod9$. \hfill\(\square\)

For a prime $p\notin\{2,3,5\}$, one also has $\tau_p\mid p-1$, since its orbit belongs to the multiplicative group of order $p-1$. The congruence of Proposition~\ref{rz-num-02-2} also holds for composites coprime to ninety. Its meaning is residual closure; irreducibility remains determined by product openings.

The numbers $7,13,77,91,143$ have decimal period six, but only the first two are irreducible. The structural relation is not that a period by itself identifies a prime, but that the product constructs components whose composition organizes returns.

\subsection{Factorization as mode occupation}
Let $\mathcal P$ be the set obtained by evaluating native irreducibles. Each mode occupation is a collection $\mathbf k=(k_p)_{p\in\mathcal P}$ of nonnegative integers with finite support. Its multiplicative evaluation is

\[
N(\mathbf k)=\prod_{p\in\mathcal P}p^{k_p}.
\]
Unique factorization identifies these occupations with the positive integers. The Hilbert spaces of occupations and of positive normal forms then have a unitary correspondence of bases. An occupation register preserves how often each irreducible occurs.

In the total-occupation-one sector, \(\sum_p k_p=1\), with basis \(\{e_p:p\in\mathcal P\}\), define

\[
H_{\mathcal P}e_p=(\log p)e_p,
\]
and, on occupations,

\[
H_{\mathrm{occ}}e_{\mathbf k}
=\left(\sum_pk_p\log p\right)e_{\mathbf k}.
\]
With maximal diagonal domains, both are self-adjoint. If $Qe_n=ne_n$, the basis identification satisfies

\[
U H_{\mathrm{occ}}U^*=\log Q.
\]
The equality follows from multiplicative composition:

\[
\sum_pk_p\log p=\log N(\mathbf k).
\]
The operator is not used to select primes retrospectively. It receives the irreducibles already constructed and converts their composition into an additive quantity. This is the precise place of the occupation representation in the article's genealogy.

\subsection{One occupation structure and two trace levels}
For $\Re s>1$, the trace on the total-occupation-one sector is

\[
P(s)=\operatorname{Tr}(e^{-sH_{\mathcal P}})=\sum_{p\in\mathcal P}p^{-s}.
\]
The trace on all occupations, by contrast, is

\[
Z(s)=\operatorname{Tr}(e^{-sH_{\mathrm{occ}}})
=\prod_{p\in\mathcal P}(1-p^{-s})^{-1}
=\sum_{n\ge1}n^{-s}.
\]
The last equality follows from unique factorization. The stated half-plane guarantees absolute convergence and permits grouping of terms. After identification of the integer realization, $Z$ agrees there with the zeta function. No table of zeta values has been used to define the modes.

\HMTresultado{Proposition}{Hierarchy of moments}{rz-num-02-3}
On $\Re s>1$, taking the logarithm branch determined by expansion of the convergent product,

\[
\log Z(s)=\sum_{k\ge1}\frac{P(ks)}{k}.
\]
\textbf{Proof.} The absolutely convergent expansion $-\log(1-z)=\sum_{k\ge1}z^k/k$, applied to each $z=p^{-s}$, gives

\[
\log Z(s)=\sum_p\sum_{k\ge1}\frac{p^{-ks}}k.
\]
Absolute convergence permits interchange of the sums. \hfill\(\square\)

The two indices preserve different functions: $p$ identifies the irreducible mode and $k$ its occupation multiplicity. Their relationship does not reduce to a list of constants: it produces a collection of moments on the same irreducible set.

\subsection{Meissel–Mertens and removal of the linear contribution}
The coordinated extension brings together the finite part

\[
B_1=\gamma+\sum_p\left(\log(1-p^{-1})+p^{-1}\right).
\]
The added term $p^{-1}$ cancels exactly the logarithm's first order. Thus,

\[
B_1=\gamma-\sum_{k\ge2}\frac{P(k)}k.
\]
Here $\gamma$ denotes the finite part $\gamma_0$ of the same functional, whose construction and bound are developed in Section~\ref{lectura:manuscrito-04}; it is not a metrological input or a decimal chosen for evaluation. The formula relates two operations on the constructed functional: extraction of the finite part and cancellation of each mode's first moment.

\HMTresultado{Proposition}{Bound on the moment tail}{rz-num-02-4}
For an integer $K\ge1$,

\[
0<\sum_{k>K}\frac{P(k)}k
\le\frac{2P(K+1)}{K+1}
\le\frac{2^{-K}}{K+1}\left(1+\frac2K\right).
\]
\textbf{Proof.} Since $k\ge K+1$,

\[
\sum_{k>K}\frac{P(k)}k
\le\frac1{K+1}\sum_p\frac{p^{-(K+1)}}{1-p^{-1}}
\le\frac{2P(K+1)}{K+1}.
\]
By comparison with all integers at least two,

\[
P(K+1)\le\sum_{n\ge2}n^{-(K+1)}
\le2^{-(K+1)}+\int_2^\infty x^{-(K+1)}\,dx.
\]
Evaluating the integral gives the second bound. \hfill\(\square\)

The inequality converts the finite sum into an interval, provided the finite evaluations and $\gamma$ themselves have controlled bounds. The length of a decimal expansion does not replace this estimate.

\subsection{Residual exclusion and the separation-two pattern product}
For the two-position pattern \(\{0,2\}\), let $\nu_p$ be the number of forbidden classes modulo $p$. At $p=2$ the two positions coincide; for $p>2$, they are two distinct classes. Local normalization relative to two independent choices produces

\[
\sigma_p=\frac{1-\nu_p/p}{(1-1/p)^2}.
\]
The factor at $p=2$ equals two, while the others form

\[
C_2=\prod_{p>2}\frac{1-2/p}{(1-1/p)^2}
=\prod_{p>2}\left(1-\frac1{(p-1)^2}\right).
\]
The pattern's complete singular series is $2C_2$. It is useful to reserve different symbols for the product over $p>2$ and for the full factor: normalization distinguishes a specific residual contribution, not a dispensable convention.

\HMTresultado{Proposition}{The same moment system with a different weight}{rz-num-02-5}
The product expansion satisfies

\[
\log C_2=-\sum_{k\ge2}\frac{2^k-2}{k}\bigl(P(k)-2^{-k}\bigr).
\]
\textbf{Proof.} For each $p>2$,

\[
\log(1-2/p)-2\log(1-1/p)
=-\sum_{k\ge2}\frac{2^k-2}{kp^k}.
\]
The sum begins at $k=2$, because the linear term cancels. For \(p\ge3\),

\[
-\log\left(1-\frac1{(p-1)^2}\right)\le\frac4{3(p-1)^2}.
\]
Comparison with the inverse-square series proves finiteness of the absolute-value sum and permits interchange of order. Removing mode $2$ produces $P(k)-2^{-k}$. \hfill\(\square\)

The same hierarchy $P(k)$ thus describes two different operations: Meissel–Mertens linear cancellation and residual exclusion of the two-position pattern. The difference lies in the weight of each moment and removal of mode $2$, not in introducing two independent prime catalogues.

For an integer \(K\ge1\), let

\[
S_K=\sum_{k=2}^K\frac{2^k-2}{k}\bigl(P(k)-2^{-k}\bigr),
\]
and the coordinated extension gives the bound

\[
0<-\log C_2-S_K
\le\frac3{K+1}\left(\frac23\right)^{K+1}\left(1+\frac3K\right).
\]
It is obtained by substituting $2^k-2\le2^k$, summing the geometric series for each $p\ge3$, bounding $(1-2/p)^{-1}\le3$ and comparing the sum over primes with the sum over integers $n\ge3$. Exponentiation gives an interval for $C_2$. Convergence of this product is a result distinct from existence of infinitely many prime pairs separated by two.

\subsection{Irreducible-mode truncation and remainder control}
\label{ix-add:corte-primos}

Occupation depth $K$ and the irreducible-mode cutoff $Y$ define distinct
approximations of the same arithmetic representation. The preceding bounds
retain every irreducible mode and truncate depth. Here the local contributions
remain complete and the modes produced by the selector are restricted to
$p\leq Y$, where $Y\geq2$ is an integer. Define
\[
 B_1(Y)=\gamma_0+\sum_{p\leq Y}
       \left(\log(1-p^{-1})+p^{-1}\right),\qquad
 C_2(Y)=\prod_{2<p\leq Y}\left(1-\frac1{(p-1)^2}\right).
\]
The empty product equals one. The finite part $\gamma_0$ is the object
denoted by $\gamma$ in the Meissel--Mertens formula; its evaluation with a
remainder is developed in Section~\ref{lectura:manuscrito-04}.

\HMTresultado{Proposition}{Enclosures from an irreducible-mode cutoff}{ix-add:prime-cut-bound}
For every integer $Y\geq2$,
\[
 B_1(Y)-\frac1{2Y}\leq B_1\leq B_1(Y),\qquad
 \frac{Y-1}{Y}C_2(Y)\leq C_2\leq C_2(Y).
\]
\textbf{Proof.}
For each $p\geq2$, the contribution subtracted from $B_1(Y)$ is
\[
 d_p=-\log(1-p^{-1})-p^{-1}
     =\sum_{k\geq2}\frac1{kp^k},\qquad
 0<d_p\leq\frac1{2p(p-1)}.
\]
The last inequality uses $1/k\leq1/2$ and the geometric sum.
Positive comparison and telescoping yield
\[
 0\leq B_1(Y)-B_1=\sum_{p>Y}d_p
 \leq\frac12\sum_{n>Y}\frac1{n(n-1)}=\frac1{2Y}.
\]
Every factor in the tail of $C_2$ lies in $(0,1)$. Including every
integer $n>Y$ decreases the product. For $M>Y$,
\[
 \prod_{n=Y+1}^{M}\left(1-\frac1{(n-1)^2}\right)
 =\prod_{n=Y+1}^{M}\frac{n(n-2)}{(n-1)^2}
 =\frac{M(Y-1)}{Y(M-1)}.
\]
Its limit is $(Y-1)/Y$. Convergence of the prime product and the
preceding comparison give the second enclosure. \hfill$\square$

As $Y$ increases, $B_1(Y)$ and $C_2(Y)$ decrease towards their respective
values. The guaranteed widths are $1/(2Y)$ and $C_2(Y)/Y$. The depth and
irreducible-mode enclosures may be intersected to refine the evaluation;
each retains its parameter and proof. The product $C_2$ corresponds to
the sector $p>2$; the local factor at mode two preserves the normalization
$2C_2$ of the complete singular series.


\subsection{Triadic realization of the functional}
The source for spectral functionals uses cycles $C_m=\mathbb Z/3^m\mathbb Z$ and their Laplacian

\[
(\Delta_mf)(j)=2f(j)-f(j+1)-f(j-1).
\]
Writing \(N_m=3^m\), let \(e_{m,n}\) be the normalized characters of the cycle and \(\widetilde n\) the representative of \(n\) between \(-(N_m-1)/2\) and \((N_m-1)/2\). The signed-orientation operator is defined by

\[
D_m e_{m,n}=\frac{N_m}{\sqrt{\pi}}
\sin\left(\frac{\pi\widetilde n}{N_m}\right)e_{m,n}.
\]
This fixes the sign, not merely the square. The realization satisfies

\[
D_m^2=\frac{3^{2m}}{4\pi}\Delta_m.
\]
Let \(P_0\) be the orthogonal projector onto constant functions on the cycle. The coordinate $\pi$ belongs to the prior generation of the HMT kernel. This normalization is part of the analytic realization and must preserve its provenance. The positive trace removes the constant mode and counts one contribution per pair of opposite orientations:

\[
\operatorname{Tr}_m^+(F(D_m))
=\frac12\operatorname{Tr}\bigl(F(D_m)P_0^\perp\bigr).
\]
With the cycle's Fourier modes, the heat trace converges to

\[
\Theta_+(x)=\sum_{n\ge1}e^{-\pi n^2x},\qquad x>0.
\]
The proof, developed in the first section of the following chapter, preserves a Gaussian majorant independent of refinement level. The Mellin transform of this limit satisfies exactly

\[
\int_0^\infty\Theta_+(x)x^{s/2-1}\,dx
=\pi^{-s/2}\Gamma(s/2)Z(s),\qquad \Re s>1.
\]
This brings together additive construction of the integer spectrum and multiplicative decomposition into irreducibles. The identity between functionals is proved; it is not inferred from proximity of two numerical evaluations.

Theta inversion and continuation of the completed function are proved in the next chapter. Finite parts, Laurent coefficients and residual classes modulo three require their specific evaluations; they do not follow merely from heat-trace convergence.

\subsubsection*{Transport defect and uniform component}
The extension on sheet folding provides an explicit representation of a one-winding defect. Let \(N=9^m\), with \(m\ge1\), and let \(T\) be a unitary operator on a Hilbert fibre \(\mathcal H\). On \(\mathbb C^N\), use the usual Euclidean inner product and the unit vector

\[
u_N=N^{-1/2}(1,\ldots,1).
\]
Define \(\iota_Nf=u_N\otimes f\), \(P_N=\iota_N\iota_N^*\) and \(Q_N=I-P_N\). Transport applies \(T\) once upon completion of the cycle:

\[
U_N(T)=\sum_{r=1}^{N-1}|r+1\rangle\langle r|\otimes I
+|1\rangle\langle N|\otimes T.
\]
It acts on the uniform state by

\[
U_N(T)\iota_N f
=\iota_N f+\frac{e_1}{\sqrt N}\otimes(T-I)f.
\]
The added term has a uniform component and an orthogonal component. The latter, with the source's normalization, is

\[
\delta_N(T)=\frac{N}{\sqrt{N-1}}Q_NU_N(T)\iota_N
=\eta_N\otimes(T-I),
\]
where

\[
\eta_N=\sqrt{\frac N{N-1}}\left(e_1-\frac{u_N}{\sqrt N}\right),
\qquad \|\eta_N\|=1,\qquad \eta_N\perp u_N.
\]
The equality is checked by substituting the preceding action and applying \(Q_N\). For two unitary transports \(T,S\) on the same fibre, one obtains

\[
\delta_N(T)^*\delta_N(S)=(I-T)^*(I-S).
\]
Commutativity is not required: the products occur in the indicated order. The same decomposition gives

\[
\delta_N(T)^*P_N\delta_N(S)=0,
\]
since \(P_N\) annihilates the factor \(\eta_N\).

These identities specify the information preserved by defect normalization and its position relative to the uniform component. They also allow the closures of the images to be transported isometrically between two cycle lengths: the Gram operator is independent of \(N\). Its domain is the constructed defect sector; it is not by itself an identification of all TPK histories or of any subsequent spectral projection.

\subsection{The completed form and its boundary identity}
Passage to the double circle uses a form on a test-function domain. Its arithmetic side brings together the prime-power, gamma and polar contributions. Writing it with two columns,

\[
\mathscr I_{\mathrm{GW}}=J_+^*J_+-J_-^*J_-,
\]
exhibits a difference of forms. The prime-index producer, orbital weights, Fourier normalization and completion must appear before it is identified with another form.

The source in Chapter 96 places the central step in an encoder $\Xi$ and an orthogonal projection $P$, through

\[
\langle\Xi f,\Xi g\rangle=\langle J_+f,J_+g\rangle,
\qquad
\langle P\Xi f,P\Xi g\rangle=\langle J_-f,J_-g\rangle.
\]
If these two identities are obtained from the declared native actions, their difference gives exactly

\[
\mathscr I_{\mathrm{GW}}(f,g)
=\langle(I-P)\Xi f,(I-P)\Xi g\rangle.
\]
The next chapter proves the explicit column decomposition and the consequence of the two moments. Application to the native encoder retains as antecedents its action on test functions, the position of its charts relative to $P$ and the proof of both moments. Orthogonality of $P$ explains the last equality but does not by itself determine the preceding two identities.

Transport telescoping, in turn, preserves the form

\[
\|x\|^2=\sum_{r=0}^{N-1}\|DA^rx\|^2+\|A^Nx\|^2,
\]
when $D^*D=I-A^*A$. If on the sheet considered $A=qV$, with $q=5/9$ and $V$ isometric, the final term equals $q^{2N}\|x\|^2$ and tends to zero. The proof of this extinction is preserved together with the identity placing state $x$ on that sheet.

\subsection{Spectral circle and memory circle}
The first circle is a coordinate of the spectral line. For $s=1/2+i\gamma$, with \(\gamma\in\mathbb R\), the transformation

\[
\tau(s)=\frac{s-1}{s}
=\frac{\gamma+i/2}{\gamma-i/2}
\]
belongs to the unit circle without the point \(1\) and is invertible through

\[
s=\frac1{1-\tau},\qquad
\gamma=\frac12\cot\frac{\theta}{2},\quad\tau=e^{i\theta}.
\]
The second circle is the character collection of integer memory:

\[
\chi_\tau(n)=\tau^n,\qquad n\in\mathbb Z.
\]
The nonadic connection increments memory by one upon completion of a winding. Intertwining relates a spectral phase to its action on that memory; it does not identify zero number $j$ with winding number $j$.

If \(C\in\mathcal B(\mathcal H)\) is the phase operator obtained in the spectral realization, the nine-position calendar acts on \(\mathbb C^9\otimes\mathcal H\) by

\[
S_{C,9}=\sum_{r=1}^8|r+1\rangle\langle r|\otimes I
+|1\rangle\langle9|\otimes C.
\]
Each vector traverses the link containing $C$ exactly once in a complete winding. Thus,

\[
S_{C,9}^9=I_9\otimes C.
\]
The identity is exact for any $C$; when $C$ is unitary, $S_{C,9}$ is also unitary. The contractive factor $5/9$ belongs to another transport readout and is distinguished from this unitary phase.

The chapter's interest does not lie in representing an arbitrary matrix on nine positions. It lies in bringing together the origin of operator $C$, its domain, the preceding encoder and compatibility with TPK cylinders. Those data fix what is transported and why this realization is relevant.

\subsection{Compatibility with the joint structure of the continuum}
The history space, two sheets, memory, incidence and complex are not reconstructed as independent objects for each section. The prior correlative assembly is developed in Section~\ref{lectura:manuscrito-05}, with its operations, realization conditions and naturality squares. The integer, depth, irreducible and spectral-height operators remain typed; coincident notation does not identify them.

Section~\ref{lectura:manuscrito-03} brings together the boundary identity with its terminal construction, test functions, terms of the completed form and convergence. Matrix approximations preserve their three matrices of form, moment and residual. Digits are compared only after outputs have been produced and frozen.

This organization preserves the requested argument: construction, operation, resulting structure and recognition. The manuscript does not change the causal kernel to facilitate a subsequent calculation, nor use a list of zeros as input to the arithmetic producer.
